\section{Introduction}
\label{sec:introduction}

% Derived from thesis sections 1.1--1.3, with scope checked against chapters 7--8.
% Internal system names and business-specific process rules are omitted.
Large language models (LLMs) can use external tools to retrieve information and invoke software functions~\cite{Schick2023Toolformer}.
In a conversational scheduling assistant, this capability must support a sequence of dependent actions: interpreting preferences, clarifying missing information, retrieving suitable options and confirming a booking.
A plausible response alone is insufficient.
The application must also retain dialogue context, validate inputs and keep the final booking decision under human control.

The Model Context Protocol (MCP) provides a common interface through which applications can access tools exposed by servers~\cite{MCP2025Architecture}.
For the scheduling task considered here, a further design decision concerns orchestration: which component selects and sequences the required operations?
The LLM can coordinate individual tools, or a backend tool can bundle several operations into a predefined sequence.
A separate planning stage introduces another design choice whose value must be assessed alongside its additional processing effort.

This paper investigates the following question: \textit{How can an MCP-based scheduling assistant be designed for human control, and how can its model and orchestration configuration be selected and evaluated?}
We examine a prototype that connects a conversational interface to existing service capabilities, maintains structured dialogue state and requires explicit user confirmation before booking.
The selection stage compares LLM orchestration with a planner, backend orchestration with a planner, and LLM orchestration without a separate planner (Figure~\ref{fig:paper-orchestration}).

The contribution is an empirical case study of controlled tool use in a bounded scheduling workflow.
It connects three aspects: a generalised design that separates model decisions from application controls, a configuration-selection study that distinguishes execution success from outcome correctness, and a subsequent assessment of the selected system through functional cases, direct prompt injection and user feedback.
The selection also considers response time, token usage and model-call cost.
Together, these results show where efficient tool execution can still fail to meet the task goal and where a working prototype needs further interaction and integration improvements.
MCP is the shared integration layer, so the study evaluates orchestration choices rather than an intrinsic advantage of the protocol.
The following sections present the background and requirements, system design, evaluation, and implications for comparable applications.
